% CC BY 4.0 License
% Copyright (c) 2026 Nanum Learn

\documentclass[a4paper,11pt]{article}

\usepackage[T1]{fontenc}
\usepackage[utf8]{inputenc}
\usepackage[english]{babel}

\usepackage{amsmath}
\usepackage{amssymb}
\usepackage{amsthm}
\usepackage{mathtools}

\usepackage[margin=1in]{geometry}
\usepackage{lastpage}
\usepackage{fancyhdr}
\fancyhead[L]{NL01CHEM01 Session 1 Problem Set}
\fancyhead[R]{Nanum Learn}
\fancyfoot[C]{}
\pagestyle{fancy}
\usepackage[skip=10pt]{parskip}
\setlength{\headheight}{13.59999pt}
\setlength{\headsep}{8pt}

\usepackage{titlesec}
\titleformat{\section}
    {\normalfont\LARGE\bfseries\color{BrickRed}}
    {\S\! \thesection}
    {0.5em}
    {}
\titleformat{\subsection}
    {\normalfont\Large\bfseries}
    {\thesubsection}
    {0.5em}
    {}

\usepackage{enumitem}
\usepackage[dvipsnames]{xcolor}


\newcommand{\easy}{(\ensuremath{\bullet \circ \circ})}
\newcommand{\medium}{(\ensuremath{\bullet \bullet \circ})}
\newcommand{\hard}{(\ensuremath{\bullet \bullet \bullet})}

\newtheoremstyle{listprobstyle}
    {\parsep}{\parsep}
    {\normalfont}
    {}
    {\bfseries}
    {}
    { }
    {{\thmname{#1}~\thmnumber{#2}.}~
        \thmnote{\textcolor{BrickRed}{\bfseries #3}}\\}
\theoremstyle{listprobstyle}
\newtheorem{problem}{Problem}



\title{\textbf{NL01CHEM01 Session 1 Problem Set}}
\date{3 November 2026}
\author{William Kang}

\begin{document}

\maketitle
\begin{center}
    \small\textit{Edited and Formatted by Enoch Yu}
\end{center}

\vspace{\baselineskip}
\begin{center}
\begin{minipage}{0.7\textwidth}
These are six chemistry practice problems written by Chemistry
director, William Kang. The problems are arranged roughly in
increasing difficulty with subjective ratings, and covers the topics
from William's presentation today. Enjoy!
\end{minipage}
\end{center}
\vspace{\baselineskip}

\section{Problems}
\begin{problem}[Element, Compound, or Mixture? \easy]
Write E, C, or M next to each.
\vspace{-0.6\baselineskip}
\begin{enumerate}[label=\alph*),itemsep=0pt,parsep=0pt,topsep=0pt]
\item Copper wire (\phantom{X})
\item Table salt (\phantom{X})
\item Well-stirred salt water (\phantom{X})
\item Room air (\phantom{X})
\item Distilled water (\phantom{X})
\end{enumerate}
\end{problem}

\begin{problem}[Pure, Uniform, or Non-Uniform? \easy]
Label each as PURE, UNIFORM, or NON-UNIFORM.
\vspace{-0.6\baselineskip}
\begin{enumerate}[label=\alph*),itemsep=0pt,parsep=0pt,topsep=0pt]
\item Spoonful of stirred sugar water:
    \underline{\phantom{XXXXXXXXXXXX}}
\item Handful of trail mix: \underline{\phantom{XXXXXXXXXXXX}}
\item Sealed distilled water: \underline{\phantom{XXXXXXXXXXXX}}
\item Rested Italian dressing: \underline{\phantom{XXXXXXXXXXXX}}
\end{enumerate}
\end{problem}

\begin{problem}[A Mystery Block \medium]
A small solid block has mass $54\,\mathrm{g}$ and volume
$20\,\mathrm{cm^3}$. (a) Compute its density. (b) Which of the
following elements is the solid block most likely made of:
aluminum ($2.7\,\mathrm{g/cm^3}$), iron ($7.9\,\mathrm{g/cm^3}$),
or copper ($8.9\,\mathrm{g/cm^3}$)?
\end{problem}

\newpage

\begin{problem}[Counting Parts of an Atom \medium]
An atom has atomic number 17 and mass number 35.
\vspace{-0.6\baselineskip}
\begin{enumerate}[label=\alph*),itemsep=0pt,parsep=0pt,topsep=0pt]
\item Which element is it?
\item Number of protons?
\item Number of neutrons?
\item Number of electrons if neutral?
\end{enumerate}
\end{problem}

\vfill

\begin{problem}[Average Atomic Mass \hard]
An element has two isotopes. Isotope M weighs $10.0\,\mathrm{u}$
and makes up 20.0\% of atoms. Isotope N weighs $11.0\,\mathrm{u}$
and makes up 80.0\%. Find the average atomic mass to one decimal
place.
\end{problem}

\vfill

\begin{problem}[Glass of Water \hard]
Use $18\,\mathrm{g/mol}$ for water and $6.02 \times 10^{23}$
particles/mol.
\vspace{-0.6\baselineskip}
\begin{enumerate}[label=\alph*),itemsep=0pt,parsep=0pt,topsep=0pt]
\item How many moles in $180\,\mathrm{g}$?
\item How many molecules total $\rightarrow$ write in scientific
notation.
\item Roughly how many per sip if a glass is six sips?
\end{enumerate}
\end{problem}

\vspace{12\baselineskip}
\end{document}


